% SGA 6, Exposición XIV: fragmento español alineado con la fuente.
% Alcance exacto: sección 4, sga6_fr_r1/src/fr.tex, líneas 15555--15725.
% SHA-256 de la fuente francesa completa:
% 0D2158B9FE3B466309A92165727A0B191422E05FBAA0DACD5806E8322C54BA1A.
% Termina inmediatamente antes de la sección 5, sin material sintético.
% Las superficies matemáticas se conservan diplomáticamente; las anomalías
% detectadas se notifican por separado para su resolución canónica.

\section*{4. \uline{Relaciones entre $K'(X)$ y el anillo de Chow $A(X)$}}

\textbf{4.1.} Supongamos, para simplificar, que $X$ sea un esquema liso y
cuasiproyectivo sobre un cuerpo $k$ (cf. n\textsuperscript{o}~8 para un caso
más general), de modo que se dispone del anillo de Chow $A(X)$ de las clases
de ciclos módulo equivalencia racional [4]. Como variantes de este anillo,
definido además bajo consideraciones sensiblemente más generales, se tienen
los anillos $\operatorname{Gr}'(X)$ (graduado asociado a la
$\lambda$-filtración de $K'(X)$) y
$\operatorname{Gr}'_{\mathrm{top}}(X)$ (graduado asociado a la filtración de
$K'(X)\simeq K_\bullet(X)$ por la codimensión de los soportes). Se tienen
teorías de las clases de Chern con valores en cada uno de estos tres anillos
(cf. [9] para el caso de $A(X)$, y este seminario para
$\operatorname{Gr}'(X)$, que se envía a
$\operatorname{Gr}'_{\mathrm{top}}(X)$). Se tiene también una teoría
cohomológica ($\ell$-ádica, digamos, donde $\ell$ es primo con la
característica de $k$) de las clases de Chern (SGA 5 VII):
\[
 c_i(\underline E)\in H^{2i}(X,\underline{\Z}_\ell(i)),
\]
(donde $\underline E$ es un fibrado vectorial sobre $X$, o también un
elemento cualquiera de $K'(X)$). Estas teorías cohomológicas están
relacionadas mediante los homomorfismos canónicos (compatibles con las
estructuras multiplicativas)
\begin{equation}\tag{4.1}
\begin{tikzcd}[column sep=large,row sep=large]
 H^{2*}(X,\underline{\Z}_\ell(*)) & {} \\
 A(X) \arrow[u,"\gamma"] \arrow[r,"\varphi"] & \operatorname{Gr}'_{\mathrm{top}}(X) \\
 {} & \operatorname{Gr}'(X) \arrow[u,"j"]
\end{tikzcd}
\end{equation}
donde se ha puesto
\[
 H^{2*}(X,\underline{\Z}_\ell(*))=\coprod_i H^{2i}(X,\underline{\Z}_\ell(i)).
\]
El homomorfismo $j$ procede del hecho de que la $\lambda$-filtración de
$K'(X)$ es más fina que la filtración topológica por la codimensión del
soporte (Exp. X), y es un isomorfismo módulo torsión (Exp. VII 4.11). El
homomorfismo $\varphi$ se define, por ejemplo, en [RRR], y asocia a la clase
del ciclo de codimensión $i$ asociado a un subesquema cerrado $Y$ de $X$ de
codimensión $i$, el elemento $\cl_X^i(\underline O_Y)$ de
$\Ftop^i(K'(X))$ módulo $\Ftop^{i+1}(K'(X))$. Es un homomorfismo
sobreyectivo, del que veremos también (4.2) que es un isomorfismo módulo
torsión. Así, \uline{módulo torsión, los tres anillos $A(X)$,
$\operatorname{Gr}'_{\mathrm{top}}(X)$ y $\operatorname{Gr}'(X)$ son
canónicamente isomorfos}; por tanto, \uline{a posteriori}, las fórmulas de
Riemann--Roch expresadas mediante uno u otro de estos anillos (tensorizados
por $\Q$) son equivalentes. El anillo graduado $\operatorname{Gr}'(X)$,
utilizado sobre todo en este seminario como sustituto del anillo de Chow,
tiene la ventaja apreciable de estar definido para todo topos localmente
anillado y, en particular, para todo esquema (regular o no). Tiene, en cambio,
la desventaja (4.6), frente a $A(X)$ y
$\operatorname{Gr}'_{\mathrm{top}}(X)$, de no poseer tal cual un carácter
covariante con respecto a los morfismos propios (incluso para esquemas lisos
cuasiproyectivos sobre un cuerpo); solo adquiere este carácter una vez
tensorizado por $\Q$.

Otra ventaja de $A(X)$ sobre sus competidores
$\operatorname{Gr}'_{\mathrm{top}}\sisi{}{(X)}$ y $\operatorname{Gr}'(X)$ es que se % Canon SGA-SRC-001448
dispone de un homomorfismo de $A(X)$ en
$H^{2*}(X,\underline{\Z}_\ell(*))$, que permite interpretar las clases de
Chern cohomológicas como imágenes de las clases de Chern con valores en
$A(X)$. En general, este homomorfismo no se factoriza por
$\operatorname{Gr}'_{\mathrm{top}}(X)$ (4.7), y es, cuando menos, dudoso que
exista en general un homomorfismo de $\operatorname{Gr}'(X)$ en
$H^{2*}(X,\underline{\Z}_\ell(*))$ que transforme clases de Chern en clases de
Chern (cf. n\textsuperscript{o}~5). Naturalmente, estas dificultades
desaparecen cuando se tensoriza por $\Q$, y se obtiene un homomorfismo
\begin{equation}\tag{4.2}
 \operatorname{Gr}'_{\mathrm{top}}(X)\longrightarrow H^{2*}(X,\underline{\Q}_\ell(*))=
 H^{2*}(X,\underline{\Z}_\ell(*))\otimes_{\underline{\Z}_\ell}\underline{\Q}_\ell
\end{equation}
que transforma clases de Chern en clases de Chern.

\textbf{4.2.} Para establecer que el homomorfismo sobreyectivo $\varphi$ de
(4.1) es un isomorfismo módulo torsión, se utiliza el homomorfismo de Chern
(deducido de la teoría de las clases de Chern [9])
\[
 K'(X)\longrightarrow \Ch(A(X))
\]
(donde el segundo miembro es el anillo de Chern de $A(X)$ definido en la Exp.
V), que, como se ve fácilmente, es compatible con la filtración topológica de
$K'(X)$ y la filtración natural de $\Ch(A(X))$, y da, al pasar a los graduados
asociados, un homomorfismo de grupos graduados
\begin{equation}\tag{4.3}
 \psi:\operatorname{Gr}'_{\mathrm{top}}(X)\longrightarrow A(X).
\end{equation}
Sentado esto, se demostrarán las fórmulas
\begin{equation}\tag{4.4}
\left\{
\begin{aligned}
 \psi\varphi(x)&=(-1)^{i-1}(i-1)!\,x\quad \text{mod. torsion}\quad \text{para }x\in A^i(X),\\
 \varphi\psi(x)&=(-1)^{i-1}(i-1)!\,x\quad \text{mod. torsion}\quad \text{para }x\in \operatorname{Gr}_{\mathrm{top}}^{\prime i}(X),
\end{aligned}
\right.
\end{equation}
de las cuales la primera implica además la segunda, gracias a que $\varphi$
es sobreyectivo. Estas fórmulas implican evidentemente que $\varphi$ y $\psi$
son isomorfismos módulo torsión y que, salvo un factor
$(-1)^{i-1}(i-1)!$ en grado $i$, son inversos el uno del otro.

Para demostrar la primera fórmula (4.4), se puede suponer evidentemente que
$x$ está definido por un subesquema cerrado irreducible $Y$ de $X$ de
codimensión $i$, en cuyo caso la fórmula equivale a
\begin{equation}\tag{4.5}
 c_i(\underline O_Y)=(-1)^{i-1}(i-1)!\,\cl(Y)\quad \text{mod. torsion},
\end{equation}
donde, en el primer miembro, $\underline O_Y$ se considera como un haz
coherente sobre $X$, que define por tanto un elemento de $K'(X)$, y, en el
segundo miembro, $\cl(Y)$ denota la clase de $Y$ en $A^i(X)$. Como $A^i(X)$
no cambia al retirar de $X$ una parte cerrada de codimensión $>i$, se reduce
al caso en que $Y$ es regular y, por tanto, liso sobre el cuerpo base $k$ si
este es perfecto. Pero, en este último caso, la fórmula (4.5) no es sino la
fórmula de Riemann--Roch para la inclusión $Y\to X$ y el elemento $1$ de
$K'(Y)$, escrita utilizando la teoría de las clases de Chern con valores en
el anillo de Chow $A(X)$. Esta fórmula se establece en [2] bajo las
condiciones consideradas aquí, mediante un método muy próximo al utilizado en
este seminario para la fórmula análoga que emplea $\operatorname{Gr}'(X)$.
Esto establece, pues, al menos para $k$ perfecto, la fórmula (4.4) y, por
consiguiente, el hecho de que $\varphi$ sea un isomorfismo módulo torsión (lo
que implica a su vez, a posteriori, que la fórmula de Riemann--Roch ``con
denominadores'', para un morfismo propio de esquemas algebraicos lisos sobre
$k$, es esencialmente la misma, tanto si se escribe utilizando $A(X)$ como en
[2], como si se utiliza $\operatorname{Gr}'(X)$ como en este seminario). El
caso en que no se hace ninguna hipótesis sobre $k$ se reduce fácilmente al caso
en que $k$ es algebraicamente cerrado, utilizando el hecho de que, para una
extensión finita $k'$ de grado $n$ de $k$, el núcleo de la aplicación natural
\[
 A(X)\xrightarrow{f^*}A(X\otimes_k k')
\]
es anulado por $n$ (como se ve utilizando el homomorfismo traza
$f_*:A(X\otimes_k k')\to A(X)$, que satisface
$f_*f^*=n\,\id$).\begingroup\renewcommand{\thefootnote}{*}\footnote{Queda,
pues, establecida gracias a P. JOUANOLOU para $X$ liso y cuasiproyectivo sobre
un cuerpo (cf. nota p.5).}\addtocounter{footnote}{-1}\endgroup

\textbf{4.3.} Parece bastante plausible que las fórmulas (4.4) sean ciertas
tal cual y no solo módulo torsión o, lo que viene a ser lo mismo, que la
fórmula (4.5) sea cierta no solo módulo torsión. Esta cuestión es un caso muy
particular de la validez en $A(X)$ de la ``fórmula de Riemann--Roch sin
denominadores'' (3.1), para una inmersión $Y\to X$ (a saber, el caso en que
$i=d$). La validez de este caso particular puede considerarse bastante
plausible. En efecto, ya se ha señalado en el n\textsuperscript{o}~3 que,
cuando $k$ es de característica cero, la fórmula de Riemann--Roch sin
denominadores es válida al menos en $\operatorname{Gr}'_{\mathrm{top}}(X)$,
lo que implica en este caso que los dos miembros de (4.5) tienen la misma
imagen en $\operatorname{Gr}'_{\mathrm{top}}(X)$ o, lo que viene a ser lo
mismo, que la segunda de las fórmulas (4.4) es válida; este resultado se
extenderá por sí mismo al caso de característica arbitraria una vez resuelta
de manera satisfactoria la cuestión planteada en el n\textsuperscript{o}~1.
Como apoyo para la fórmula (4.5) en el propio $A(X)$, señalemos el hecho de que
es claro a priori que la restricción del primer miembro a $X-Y$ es nula; por
tanto, suponiendo $Y$ irreducible, el primer miembro debe ser a priori un
múltiplo entero de $\cl(Y)$:
\[
 c_i(\sisi{Y}{\mathcal O_Y})=n(Y)\cl(Y). % Canon SGA-SRC-001449
\]
Esto demuestra ya la fórmula exacta (4.5) cuando $\cl(Y)$ no es un elemento
de torsión de $A^i(X)$. Es inmediato que este es el caso si $X$ es proyectivo
(y no solo cuasiproyectivo) sobre el cuerpo $k$, observando que, si
$\xi\in A^1(X)$ es una polarización de $X$, entonces
$\deg(Y\cdot\xi^{\sisi{i}{\dim Y}})>0$; esto demuestra la fórmula deseada en este caso. Se % Canon SGA-SRC-001450
deduce que esta fórmula sigue siendo válida siempre que $X$ sea isomorfo a un
subesquema abierto de un esquema proyectivo liso sobre $k$, lo que ocurrirá
cuando se disponga de una solución adecuada del problema de resolución de
singularidades. Así sucede, en particular, si $k$ es de característica cero,
gracias a Hironaka [11].

La validez de la fórmula exacta (4.5), y por tanto de (4.4), implicaría que el
núcleo de
\[
 \varphi^i:A^i(X)\longrightarrow \operatorname{Gr}_{\mathrm{top}}^{\prime i}(X)
\]
es anulado por $(i-1)!$ y, en particular, que
$\varphi^2:A^2(X)\to\operatorname{Gr}_{\mathrm{top}}^{\prime 2}(X)$ es un
isomorfismo. (Que $\varphi^1$ sea un isomorfismo se deduce ya inmediatamente
de X 5.3.2). Salvo error, el hecho de que $\varphi^2$ sea un isomorfismo no
está demostrado en la actualidad, ni siquiera en el caso particular en que
$\dim X=3$. En cualquier caso, fue precipitadamente como el redactor había
afirmado la validez de (4.4) y (4.5) en [9,
n\textsuperscript{o}~4, 3\textsuperscript{o}].

\textbf{4.4.} En la actualidad, la principal ventaja de
$\operatorname{Gr}'_{\mathrm{top}}(X)$ (o de $\operatorname{Gr}'(X)$) sobre
el anillo de Chow $A(X)$ es que se presta mejor a los cálculos explícitos de
ciertas fórmulas de intersección entre ciclos, en el caso de intersecciones
``excedentarias'', es decir, de dimensión demasiado grande. La definición de
la clase de intersección, mediante el lema de Chow que permite mover
suficientemente un ciclo dentro de su clase, se presta, en efecto, con
dificultad al cálculo. Un ejemplo notable a este respecto es la fórmula
fundamental siguiente [Exp. VII 2.7], para una inmersión cerrada $f:Y\to X$
de esquemas cuasiproyectivos lisos sobre $k$:
\begin{equation}\tag{4.6}
 f^*f_*(y)=y\,c_d(\underline N),
\end{equation}
fórmula válida en $\operatorname{Gr}'_{\mathrm{top}}(Y)$, donde $y$ está en
$\operatorname{Gr}'_{\mathrm{top}}(Y)$ y $\underline N$ es el haz conormal de
$Y$ en $X$, supuesto de rango $d$ en todas partes. En loc. cit., esta fórmula
se deduce de un cálculo muy sencillo de haces Tor. Gracias a lo visto en 4.2,
la fórmula anterior implica la misma fórmula en $A(Y)$, pero solo módulo
torsión. Es bastante extraordinario que la fórmula (4.6) en el propio $A(Y)$
no esté demostrada en la actualidad\begingroup\renewcommand{\thefootnote}{*}\footnote{Cf.
nota de la página 23.}\addtocounter{footnote}{-1}\endgroup, ni siquiera el
caso particular obtenido tomando $y=1$, que da la autointersección de $Y$:
\begin{equation}\tag{4.7}
 f^*(\cl(Y))=c_d(\underline N),
\end{equation}
¡incluso si $k$ es el cuerpo de los complejos! Se tropieza con un problema
análogo cuando se intenta calcular el anillo de Chow de la explosión
$X'$ de $X$ a lo largo de $Y$, imitando el método seguido para calcular
$\operatorname{Gr}'(X')_\Q$ [Exp. VII 4.8], y que funciona también (sin
tensorizar por $\Q$) para calcular $\operatorname{Gr}'_{\mathrm{top}}(X')$.
La fórmula clave para esta determinación (con las notaciones de loc. cit.) es
\begin{equation}\tag{4.8}
 f^*(i_*(y))=j_*\bigl(g^*(y)c_{d-1}(\underline{\check P})\bigr),
\end{equation}
fórmula establecida actualmente en $\operatorname{Gr}'_{\mathrm{top}}(X')$
(para $y\in\operatorname{Gr}'_{\mathrm{top}}(Y)$) mediante un cálculo de Tor
(Exp. VII 3.4), y cuya validez en $A(X')$ (para $y\in A(Y)$) quedaría por
establecer.

Señalemos que ni siquiera los análogos $\ell$-ádicos de las fórmulas (4.6),
(4.7) y (4.8) están demostrados en la actualidad; cuando $y$ es la clase de
cohomología de un ciclo algebraico, (4.6) y (4.8) son ciertas, en cualquier
caso, módulo torsión.

\textbf{4.5.} En esta sección y en las siguientes se darán ejemplos en los
que, en el diagrama (4.1), el homomorfismo $j$, respectivamente el
homomorfismo $\varphi$, no es un isomorfismo. Evidentemente, para un $X$ dado,
viene a ser lo mismo decir que
\[
 j:\operatorname{Gr}'(X)\longrightarrow \operatorname{Gr}'_{\mathrm{top}}(X)
\]
es un isomorfismo, un monomorfismo o un epimorfismo. Además, la imagen de $j$
es el subanillo de $\operatorname{Gr}'_{\mathrm{top}}(X)$ generado por las
clases de Chern de los haces localmente libres sobre $X$. En el caso en que
$K'(X)$ está generado como álgebra sobre $\Z$ por las clases de los haces
invertibles sobre $X$, se ve de inmediato que el subanillo considerado de
$\operatorname{Gr}'_{\mathrm{top}}(X)$ es también el subanillo generado por
$\operatorname{Gr}'^{1}_{\mathrm{top}}(X)$, la componente de grado $1$ de
$\operatorname{Gr}'_{\mathrm{top}}(X)$.

Tomemos ahora como $X$ una variedad de banderas
\[
 X=G/B,
\]
donde $G$ es un grupo algebraico semisimple simplemente conexo sobre el
cuerpo algebraicamente cerrado $k$, y $B$ un subgrupo de Borel [5]. Cuando
$k=\C$, la descomposición celular de Bruhat de $X$ [5, Exp. 13] implica que
el homomorfismo canónico
\[
 (*)\qquad K'(X)\longrightarrow K(X(\C))
\]
es un isomorfismo, al igual, por otra parte, que el homomorfismo canónico
\[
 (**)\qquad A(X)\longrightarrow H^*(X(\C),\Z),
\]
donde $X(\C)$ es el espacio compacto de los puntos complejos de $X$, provisto
de su topología usual deducida de la topología de $\C$. Por otra parte, L.
HODGKIN [12] demostró que $K(X(\C))$ está ciertamente generado por las clases
de los haces invertibles sobre $X$. Por tanto, lo mismo ocurre con el propio
$K'(X)$, y no es difícil deducir de ello, mediante una técnica de
especialización que utiliza los grupos semisimples desplegados sobre $\Z$
(SGA 3), que el mismo resultado sigue siendo válido para todo cuerpo $k$
algebraicamente cerrado. Además, la descomposición celular de $X$ (o la
biyectividad de $(**)$ en el caso $k=\C$) implica fácilmente que $A(X)$ no
tiene torsión; por tanto, gracias a 4.2, el homomorfismo
\[
 \varphi:A(X)\longrightarrow\operatorname{Gr}'_{\mathrm{top}}(X)
\]
es un isomorfismo. De estos hechos se deduce que $j$ es un isomorfismo si y
solo si $A(X)$ está generado por sus elementos de grado $1$ o, también, si
$k=\C$, si y solo si la cohomología entera de $X(\C)$ está generada por los
elementos de grado $2$, es decir, si el espacio clasificante $B_{G(\C)}$ no
tiene torsión. Se sabe, por otra parte [15, p. 217 th. B
(1)$\Leftrightarrow$(2), y 2.5], que esto sucede si y solo si $G$ es ``sin
torsión'', o también si $G$ es isomorfo a un producto de grupos
$\operatorname{S\ell}(n)$ y $\operatorname{Sp}(m)$. Basta, pues, tomar como $G$
un grupo $\operatorname{\sisi{SO}{Spin}}(n)$ con $n\geq 7$, o un grupo simple excepcional, % Canon SGA-SRC-001451
para obtener un ejemplo en que $j$ no es un isomorfismo.

\textbf{4.6.} Obsérvese que, en este ejemplo, no cabe duda de que el anillo
``bueno'' es el anillo sin torsión
\[
 A(X)\simeq\operatorname{Gr}'_{\mathrm{top}}(X),
\]
que tiene como base las celdas de la descomposición de Bruhat, anillo también
isomorfo a la cohomología entera de $X(\C)$ si $k=\C$; en cambio, el anillo
$\operatorname{Gr}'(X)$ aparece como un anillo con una torsión de naturaleza
bastante patológica. Este ejemplo confirma, pues, la concepción según la cual
el invariante $\operatorname{Gr}'(X)$ apenas es interesante salvo módulo
torsión, es decir, tensorizado por $\Q$. Al mismo tiempo se ve, mediante este
ejemplo, que solo una vez tensorizado por $\Q$ posee
$\operatorname{Gr}'(X)$ un carácter covariante con respecto a los morfismos
propios, incluso
limitándose a las inmersiones cerradas de esquemas cuasiproyectivos y lisos.
En efecto, si $X$ es tal que $j$ no sea un isomorfismo, existe en $X$ un
subesquema cerrado integral $Y$ tal que
\[
 \sisi{\cl^i(Y)}{\cl'(Y)}\in K'(X) % Canon SGA-SRC-001452
\]
no está en $\Filt^i(X)$, donde $i=\operatorname{codim}(Y,X)$. Se puede tomar
$Y$ de dimensión mínima, de modo que
\[
 \Filt^{i+1}(X)=\Filt_{\mathrm{top}}^{i+1}(X).
\]
Si $Z$ es el conjunto singular de $Y$, y $X'=X-Z$, $Y'=Y-Z$, se concluye que
\[
 \cl'(Y')\in K'(X')
\]
no está en $\Filt^i(X')$ (teniendo en cuenta que
$K'(X)\to K'(X')$ es sobreyectivo y compatible con los aumentos, y que su
núcleo es
\[
 \operatorname{Im}\bigl(\sisi{K(Z)\longrightarrow K(X)}{K'(Z)\longrightarrow K'(X)}\bigr)\subset % Canon SGA-SRC-001453
 \Filt_{\mathrm{top}}^{i+1}(X)=\Filt^{i+1}(X)).
\]
Consideremos entonces la inclusión
\[
 f:Y'\longrightarrow X',
\]
que es una inmersión cerrada de codimensión $i$ de esquemas
cuasiproyectivos y lisos sobre $k$, tal que
\[
 f_*:K'(Y')\longrightarrow K'(X')
\]
no es compatible con las $\lambda$-filtraciones (salvo desplazamiento por
$i$), puesto que $f_*(1)$ no está en $\Filt^i(X')$.

\textbf{4.7.} Conservando las notaciones de 4.5, observemos ahora que de (Exp.
IX) se deduce que el teorema de HODGKIN citado equivale a que se tenga un
isomorfismo
\[
 K'(G)\simeq\Z,\qquad\text{por tanto}\qquad \operatorname{Gr}'_{\mathrm{top}}(G)\simeq\Z,
\]
siendo el homomorfismo $K'(G/B)\to K'(G)$ a priori sobreyectivo y teniendo
como núcleo el ideal generado por los $\cl'(\underline L)-1$, donde
$\underline L$ recorre los módulos invertibles sobre $X=G/B$ asociados a las
representaciones $B\to\Gm$ (que dan, de hecho, todos los módulos invertibles
sobre $X$ [5, Exp. 15]). Por otra parte, un razonamiento análogo muestra que
\[
 A(G/B)\longrightarrow A(G)
\]
es también sobreyectivo y que su núcleo es el ideal generado por los elementos
de $A^1(G/B)$ definidos por las representaciones $B\to\Gm$ y, por tanto, de
hecho por todo $A^1(G/B)$.\begingroup\renewcommand{\thefootnote}{*}\footnote{Cf.
IX 3.4, 2.2.}\addtocounter{footnote}{-1}\endgroup Esto muestra que $A(G)$ se
reduce a $A^0(\sisi{X}{G})=\Z$ si y solo si ``$G$ no tiene torsión'', es decir, si $G$ % Canon SGA-SRC-001454
es isomorfo a un producto de grupos $\operatorname{S\ell}(n)$ y
$\operatorname{Sp}(m)$. Por tanto, en este caso, y solo en este caso, el
homomorfismo
\[
 \varphi:A(G)\longrightarrow\operatorname{Gr}'_{\mathrm{top}}(G)
\]
es un isomorfismo. Por tanto, si $G$ tiene torsión, el esquema subyacente
proporciona un ejemplo de un $\varphi$ no biyectivo.

\textbf{4.8.} De igual modo se obtiene un ejemplo en que el homomorfismo
$\gamma$ de (4.1)
\begin{equation}\tag{4.9}
 \gamma:A(G)\longrightarrow H^{2*}(G,\underline{\Z}_\ell(*))
\end{equation}
no se factoriza por $\operatorname{Gr}'_{\mathrm{top}}(G)$ o, lo que viene a
ser lo mismo aquí, un ejemplo en que este homomorfismo no es nulo en grados
$>0$. Limitándonos (como es lícito gracias a teoremas de especialización
adecuados) al caso en que $k=\C$, se reduce al estudio del homomorfismo
\begin{equation}\tag{4.10}
 A(G)\longrightarrow H^*(G(\C),\Z),
\end{equation}
deducido, por otra parte, pasando al cociente del homomorfismo canónico
\[
 A(G/B)\simeq H^*(G/B(\C),\Z)\longrightarrow H^*(G(\C),\Z).
\]
Es tentador conjeturar que el homomorfismo (4.10) es de hecho inyectivo, pues
se sabe, por otra parte, que los números primos de torsión para $G(\C)$ son
exactamente los que dividen el orden del grupo finito
\[
 A^+(G)=\bigoplus_{i>0}A^i(G).
\]
Si así fuera, en cuanto $G$ tuviera torsión, el homomorfismo (4.10) no se
factorizaría por $\varphi$; por tanto, para un $\ell$ adecuado (más
precisamente, para $\ell$ que divida el orden de $A^+(G)$, es decir, $\ell$
coeficiente de torsión de $G$), lo mismo ocurriría con (4.9). Al menos se puede
afirmar que, en el caso del grupo $G=\underline G_2$ sobre el cuerpo $\C$, el
homomorfismo (4.10) no es nulo en grado $>0$; más precisamente, existe un
elemento de $A^3(G)$ cuya imagen en $H^6(G(\C),\Z)$ es de orden $2$; se
deduce, pues, que el homomorfismo (4.9) no se factoriza por $\varphi$ cuando
$\ell=2$. El hecho utilizado relativo a $\underline G_2$, que equivale a decir que existe
en $H^6(G/B(\C),\Z)$ un elemento cuya imagen en $H^6(G(\C),\Z)$ no es nula,
se encuentra en [3], mediante el diccionario conocido entre el lenguaje de los
grupos algebraicos semisimples complejos y el de los grupos de Lie compactos;
me fue comunicado amablemente por J.P. SERRE.
